%=============================================================
%  pages/exp1.tex — Experiment 2
%  
%=============================================================

\experiment{3}{%
    Construct a frequency distribution table from a given dataset by grouping the data into class internvals. Represent the data using graphs such as histograms or bar charts.
}



% frequency distribution table
% histogram idk of whose
% bar chart grade distribution
% ogive chart, idk of whose

% obsertions to make
% which class internval has the highest frequnecy
% type of skewness
% 

% ── 1. Dataset Info ───────────────────────────────────────────
\subsection{Frequency Distribution Table}
We can group the dataset into class intervals by using `cut()` function of pandas.

\begin{lstlisting}[style=python]
import pandas as pd 
import matplotlib.pyplot as plt 
import numpy as np

df_student = pd.read_csv("dataset.csv", delimiter=',', header='infer')

sturges = round(1+np.log2(df_student.Marks.count()))

freq = pd.cut(df_student["Marks"], bins=sturges, include_lowest=True).value_counts().sort_index()

freq_dis = pd.DataFrame({
    'Class Interval': freq.index,
    'Frequency': freq.values,
    'Relative frequency': (freq.values / len(df_student)).round(2),
    'Cumulative Frequency': freq.values.cumsum()
})

print(freq_table)
\end{lstlisting}


\textbf{Output:}\\
\begin{tabular}{ccccc}
    \toprule
    & \textbf{Class Interval} & \textbf{Frequency} & \textbf{Relative Frequency} & \textbf{Cumulative Frequency} \\
    \midrule
    0 & 35-50  & 5 & 0.17 & 5 \\
    1 & 51-60  & 6 & 0.2 & 11 \\
    2 & 61-70 & 1 & 0.03 & 12 \\
    3 & 71-80  & 9 & 0.3 & 21 \\
    4 & 81-90  & 4 & 0.13 & 25 \\
    5 & 91-100  & 5 & 0.17 & 30 \\

    \bottomrule
\end{tabular}


% \subsection{Classification of Dataset}
% \begin{tabularx}{\linewidth}{ccccX}
%     \toprule
%     & \textbf{Variable} & \textbf{Data Type} & \textbf{Measurement Scale} & \textbf{Justification} \\
%     \midrule
%     1 & Name  & Categorial & Nominal & no mathematical relation or order \\
%     2 & Gender & Categorial & Nominal & Two categories, no relation \\
%     3 & Grade  & Categorial & Ordinal & A > B > C > D > E, intervals are unequal \\
%     4 & Marks  & Numerical & Ratio & 0 is nil, 90 marks is twice 45 marks \\
%     5 & StudyHours  & Numerical & Ratio & 0 means nil \\
%     6 & Attendance  & Numerical & Ratio & 0\% means nil \\

%     \bottomrule
% \end{tabularx}
